\documentclass[11pt,a4paper]{article}
\usepackage[T1]{fontenc}
\usepackage[utf8]{inputenc}
\usepackage{lmodern,amsmath,amssymb,amsthm,booktabs,microtype,xcolor,fancyhdr}
\usepackage[margin=26mm,bottom=23mm,headheight=15pt]{geometry}
\usepackage[colorlinks=true,linkcolor=blue!45!black,citecolor=blue!45!black,urlcolor=blue!45!black]{hyperref}
\providecommand{\Esp}{0}
\ifnum\Esp=1
\usepackage[spanish,es-nodecimaldot]{babel}
\newcommand{\B}[2]{#2}
\else
\newcommand{\B}[2]{#1}
\fi
\newcommand{\e}{\mathrm e}
\newcommand{\one}{\mathbf 1}
\newcommand{\Sg}{\mathfrak S}
\newcommand{\pos}[1]{(#1)_+}
\newcommand{\Z}{\mathbb Z}
\pagestyle{fancy}\fancyhf{}
\fancyhead[L]{\small\B{RESEARCH PREPRINT}{PREPUBLICACIÓN DE INVESTIGACIÓN}}
\fancyhead[R]{\small\B{November 2026}{Noviembre de 2026}}
\fancyfoot[L]{\small\B{Reflected prime pairs}{Parejas de primos reflejadas}}
\fancyfoot[R]{\small\thepage}
\renewcommand{\headrulewidth}{0.35pt}
\newcommand{\papertitle}{\B{Reflected prime pairs on $3+2a$: exact sieving, a three-case conjecture and a Type II obstruction}{Parejas de primos reflejadas sobre $3+2a$: criba exacta, conjetura de tres casos y obstáculo de Tipo II}}
\hypersetup{pdftitle={\B{Reflected prime pairs: exact sieving and a Type II obstruction}{Parejas de primos reflejadas: criba exacta y obstáculo de Tipo II}},pdfauthor={Tomás Sánchez Expósito}}
\setlength{\emergencystretch}{2em}
\begin{document}
\begin{center}
{\small\scshape\B{Working paper $\cdot$ Version 1.2}{Documento de trabajo $\cdot$ Versión 1.2}}\\[0.7em]
{\LARGE\bfseries\papertitle\par}\vspace{0.8em}
{\large Tomás Sánchez Expósito}\\[0.2em]
\B{Independent researcher}{Investigador independiente}\\[0.25em]
\B{22 November 2026}{22 de noviembre de 2026}
\end{center}
\begin{abstract}
\B{We connect an additive prime sieve and a batch-residue test for twin-prime candidates to the exact counting of distinct prime pairs with a prescribed even sum. The author's three-case conjecture is equivalent to distinct-prime coverage on the set where both the predecessor and the midpoint are composite. Strong Goldbach additionally requires coverage where the predecessor is prime and the midpoint is composite. We derive a square-start reflected sieve and an equivalent complete-batch GCD test with explicit small-prime exceptions. Three deterministic computations agree on every even target from 4 to 7930, transferring the published 1000-prime prefix to a two-prime counting experiment. A separate exact witness audit through $8\cdot10^6$ finds a distinct-prime pair at every even target except 4 and 6. We consolidate the project’s variance, Fourier and deletion obstructions and give the details of the Vaughan Type I reduction, the singular-series evaluation of its kernel and the centering of the Type II term. A target-uniform signed lower bound for that term would suffice for eventual coverage, but is not established. We also derive, from an external exceptional-zero formula, a conditional restriction on real Dirichlet zeros from an uncut two-sided estimate on the residual set. An exact counterexample tests an antidiagonal symmetry asserted in a recent preprint. The inherited analytic derivations await independent human review. Neither universal conjecture is proved.}{Conectamos una criba aditiva de primos y un criterio por lotes de residuos para candidatos a primos gemelos con el recuento exacto de parejas de primos distintos de suma par prefijada. La conjetura de tres casos del autor equivale a la cobertura por primos distintos cuando el predecesor y el punto medio son compuestos. Goldbach fuerte exige además cobertura cuando el predecesor es primo y el punto medio es compuesto. Deducimos una criba reflejada con arranque en el cuadrado y un criterio equivalente de máximo común divisor con lote completo y excepciones explícitas para los primos pequeños. Tres cálculos deterministas coinciden en cada objetivo par de 4 a 7930 y trasladan el prefijo publicado de mil primos a un experimento de recuento de parejas. Una comprobación exacta separada hasta $8\cdot10^6$ encuentra una pareja de primos distintos para todos los pares salvo 4 y 6. Integramos los obstáculos de varianza, Fourier y borrado del proyecto y detallamos la reducción de Tipo I de Vaughan, la evaluación de su núcleo mediante la serie singular y el centrado del término de Tipo II. Una cota inferior con signo, uniforme en el objetivo para ese término, bastaría para una cobertura eventual, pero no se obtiene. También deducimos, a partir de una fórmula externa para ceros excepcionales, una restricción condicional sobre ceros reales de Dirichlet bajo una estimación bilateral sin corte en el conjunto residual. Un contraejemplo exacto examina una simetría de antidiagonales afirmada en un preprint reciente. Las derivaciones analíticas heredadas están pendientes de revisión humana independiente. No se demuestra ninguna de las dos conjeturas universales.}
\end{abstract}
\noindent\textbf{\B{Keywords}{Palabras clave}:} \B{binary Goldbach, additive sieve, reflected prime pairs, singular series, Vaughan's identity, parity barrier}{Goldbach binario, criba aditiva, parejas de primos reflejadas, serie singular, identidad de Vaughan, barrera de paridad}.\\
\textbf{MSC 2020:} 11P32, 11N35, 11N36, 11Y11.

\smallskip\noindent\textbf{\B{AI disclosure}{Declaración sobre IA}.}
\B{OpenAI ChatGPT assisted with source comparison, mathematical derivation and exposition, experimental code, and the English and Spanish manuscripts. The new arguments are presented for scrutiny; neither AI assistance nor finite computations constitute peer review or formal certification.}{OpenAI ChatGPT se utilizó como apoyo en la comparación de fuentes, la derivación y exposición matemática, el código experimental y la preparación de los textos en inglés y español. Los argumentos se presentan para su examen; ni la asistencia de IA ni los cálculos finitos constituyen revisión por pares o certificación formal.}

\section{\B{Scope and relationship to prior work}{Alcance y relación con los trabajos previos}}
\B{This paper consolidates the Goldbach research line (project stages 56--76) rather than reproducing its internal reports. The two earlier preprints have sharply different uses. The first \cite{additive}, §§3--7 and Appendix B, formalizes an operational sieve: a newly discovered odd prime $\ell$ inserts a marker at $\ell^2$, then each encounter advances it by $2\ell$. Its Lean 4 result concerns the specified single-prime stream, including the published list of 1000 odd primes through 7927. The second \cite{twins}, §§4--8, groups forbidden residues for two endpoints into a single GCD against a complete odd-prime batch. The shifts there have a fixed gap of 2, and the final Lean statements concern the specified structured base. Here it is the sum, not the gap, that is fixed. The use of the sieve output and the batch criterion is exact, but the new reflected construction is not covered by either Lean proof.}{Este artículo consolida la línea de Goldbach (etapas 56--76 del proyecto), sin reproducir sus informes internos. Las dos prepublicaciones anteriores tienen usos distintos. La primera \cite{additive}, §§3--7 y apéndice B, formaliza una criba operativa: un primo impar recién descubierto $\ell$ inserta un marcador en $\ell^2$ y, al alcanzarlo, este avanza $2\ell$. El resultado en Lean 4 concierne al generador individual especificado, incluida la lista de mil primos impares hasta 7927. La segunda \cite{twins}, §§4--8, agrupa las clases prohibidas de dos extremos en un solo MCD frente a un lote completo de primos impares. Allí los desplazamientos tienen diferencia fija de 2 y los resultados finales en Lean se refieren a la base estructurada especificada. Aquí se fija la suma, no la diferencia. El uso de la salida de la criba y del criterio por lotes es exacto, pero ninguna de las dos pruebas en Lean abarca la nueva construcción reflejada.}

\B{Our contribution relative to those papers is the explicit reflected specialization, including the square-start intervals and the restoration of small endpoints, its finite cross-check at the published prefix, and a synthesis of the analytical dependencies. The underlying Goldbach sieve on $k(N-k)$ is classical; see Ford \cite[§1.2]{ford}. We claim no priority for the local residue calculation, no speedup, and no new signed prime-correlation estimate. The first missing theorem is a signed pointwise bound for the actual prime correlation (Section~\ref{sec:stop}). The documentary and finite comparison is at levels L0--L1; the inherited positivity implication is L2. No direct proof stage is reopened.}{La aportación respecto a esos artículos es la especialización reflejada explícita, con tramos de arranque en el cuadrado y recuperación de extremos pequeños, su contraste finito sobre el prefijo publicado y la síntesis de las dependencias analíticas. La criba de Goldbach sobre $k(N-k)$ es clásica; véase Ford \cite[§1.2]{ford}. No se reivindica prioridad sobre el cálculo de residuos locales, aceleración del algoritmo ni una nueva estimación con signo de correlaciones de primos. El primer teorema que falta es una cota puntual con signo para la correlación de los primos auténticos (sección~\ref{sec:stop}). La comparación documental y finita corresponde a L0--L1 y la implicación de positividad heredada a L2. No se reabre un ataque directo.}

\B{Background. Oliveira e Silva, Herzog and Pardi verified the binary statement through $4\cdot10^{18}$ \cite{verify}; that is a published finite baseline, not a claim about the latest possible computational record. Chen's theorem gives a prime plus a number with at most two prime factors for sufficiently large even integers \cite{chen}. Pintz's $X^{0.72}$ bound is an exceptional-set estimate for large $X$ \cite{pintz}, not a bound of zero, and Helfgott's theorem concerns odd integers greater than 5 and three summands \cite{helfgott}. None supplies the signed bound for every target needed below.}{Contexto. Oliveira e Silva, Herzog y Pardi verificaron el enunciado binario hasta $4\cdot10^{18}$ \cite{verify}; es una referencia finita publicada, no una afirmación sobre el último récord computacional posible. El teorema de Chen obtiene un primo más un número con a lo sumo dos factores primos para pares suficientemente grandes \cite{chen}. La cota $X^{0.72}$ de Pintz es una estimación del conjunto excepcional para $X$ grande \cite{pintz}, no una cota de cero, y el teorema de Helfgott concierne a impares mayores que 5 y tres sumandos \cite{helfgott}. Ninguno aporta la cota con signo que se necesita aquí para cada objetivo.}

\section{\B{The author's statement and its exact logical scope}{Enunciado del autor y alcance lógico exacto}}
\B{For each even integer $N\ge4$ let $R(N)$ count unordered distinct-prime pairs, represented with the smaller prime first. The author's conjecture is that at least one of the following cases holds for each such $N$:}{Para cada entero par $N\ge4$, sea $R(N)$ el número de parejas no ordenadas de primos distintos, representadas con el primo menor primero. La conjetura del autor sostiene que para cada $N$ se cumple al menos uno de los casos siguientes:}
\[
 N=p+1,\qquad N=2p,\qquad N=p+q\quad(p<q),
 \qquad p,q\ \B{\text{prime}}{\text{primos}}.
\]
\B{The cases can overlap. In the third case the reflected-doubles identity $N-2p=2q-N=q-p$ follows from the sum, and equality of the two outer expressions conversely forces $p+q=N$. In particular, $q=N-p$ must itself be prime; the equality does not test its primality. In particular $1$ is not prime; the first case need not itself give a representation by two primes.}{Los casos pueden solaparse. En el tercero, la identidad de los dobles reflejados $N-2p=2q-N=q-p$ se sigue de la suma y, a la inversa, la igualdad entre las dos expresiones exteriores fuerza $p+q=N$. En particular, $q=N-p$ ha de ser primo por derecho propio: la igualdad no certifica su primalidad. En particular, $1$ no es primo: el primer caso no proporciona por sí mismo una representación por dos primos.}
\begin{align*}
 H&=\{N\ge4:N\ \B{\text{even},\ N-1\text{ and }N/2\text{ composite}}{\text{par},\ N-1\text{ y }N/2\text{ compuestos}}\},\\
 H_1&=\{N\ge8:N\ \B{\text{even},\ N-1\text{ prime},\ N/2\text{ composite}}{\text{par},\ N-1\text{ primo},\ N/2\text{ compuesto}}\},\\
 D&=\{N\ge4:N\ \B{\text{even},\ N/2\text{ composite}}{\text{par},\ N/2\text{ compuesto}}\}=H\mathbin{\dot\cup} H_1.
\end{align*}
\B{Proof of the two reductions. For even $N\ge4$, both $N-1$ and $N/2$ are greater than 1. Outside $H$ one of the first two author cases holds; inside $H$ neither holds. Thus the universal three-case conjecture is equivalent to $R(N)>0$ for every $N\in H$. A representation with equal primes exists exactly when $N/2$ is prime; therefore strong Goldbach is equivalent to $R(N)>0$ for every $N\in D$. A sum involving 2 at an even target must be $2+2$. This handles $N=4$ and explains why no distinct pair is lost by using only odd endpoints. Goldbach implies the author’s conjecture; a converse based only on that conjecture is not established. For example at $N=8\in H_1$ the case $7+1$ does not supply the actual prime pair $3+5$.}{Prueba de las dos reducciones. Para $N\ge4$ par, tanto $N-1$ como $N/2$ son mayores que 1. Fuera de $H$ vale uno de los dos primeros casos del autor; dentro de $H$ no vale ninguno. La conjetura universal de tres casos equivale, por tanto, a $R(N)>0$ para todo $N\in H$. Existe una representación con primos iguales exactamente cuando $N/2$ es primo; así, Goldbach fuerte equivale a $R(N)>0$ para todo $N\in D$. Una suma de dos primos con el 2 y objetivo par ha de ser $2+2$. Esto resuelve $N=4$ y explica por qué no se pierde ninguna pareja distinta al usar solo extremos impares. Goldbach implica la conjetura del autor; no se ha establecido el recíproco partiendo solo de esta. Por ejemplo, en $N=8\in H_1$ el caso $7+1$ no proporciona la pareja prima real $3+5$.}


\paragraph{\B{Density and exceptional sets.}{Densidad y conjuntos excepcionales.}}
\B{Of the even targets from $4$ to $X$, at most $\pi(X-1)+\pi(X/2)$ lie outside $H$, by counting prime predecessors and prime midpoints. Hence $H$ has relative density one by the prime number theorem. This neither makes the two conjectures equivalent nor makes Type II control necessary for a different proof. If $E(X)$ counts Goldbach exceptions through $X$, then $E_H(X):=\#\{N\in H:N\le X,\ R(N)=0\}\le E(X)$, since there is no prime diagonal on $H$. In particular $E_H(4\cdot10^{18})=0$ by the published verification \cite{verify}. Zhao's January 2026 preprint revision \cite{zhao} gives an unconditional $O(X^{7/10})$ bound with an ineffective constant, improving the cited $0.72$ benchmark. Pintz and Althöfer report $E(X)\ll_\epsilon X^{7/12+\epsilon}$ under a common fixed zero-free half-plane $\Re s>\theta_0$, $\theta_0<1$, using Pintz's explicit formula and the exponential-sum estimate of Maynard, Pandey and Radziwiłł \cite{pa,mpr}. We do not use this hypothesis in any proof. No positive power bound rules out a prescribed exception.}{Entre los pares de $4$ a $X$, a lo sumo $\pi(X-1)+\pi(X/2)$ quedan fuera de $H$, al contar predecesores y puntos medios primos. Por el teorema de los números primos, $H$ tiene densidad relativa uno. Eso no hace equivalentes las dos conjeturas ni obliga a controlar Tipo II en una vía de prueba distinta. Si $E(X)$ cuenta las excepciones a Goldbach hasta $X$, entonces $E_H(X):=\#\{N\in H:N\le X,\ R(N)=0\}\le E(X)$, pues en $H$ no hay diagonal prima. En particular, $E_H(4\cdot10^{18})=0$ por la verificación publicada \cite{verify}. La revisión de enero de 2026 del preprint de Zhao \cite{zhao} da la cota incondicional $O(X^{7/10})$, con constante inefectiva, que mejora la referencia citada de exponente $0.72$. Pintz y Althöfer anuncian $E(X)\ll_\epsilon X^{7/12+\epsilon}$ bajo un semiplano común fijo libre de ceros $\Re s>\theta_0$, $\theta_0<1$, usando la fórmula explícita de Pintz y la estimación de sumas exponenciales de Maynard, Pandey y Radziwiłł \cite{pa,mpr}. No se usa esa hipótesis en las pruebas. Ninguna cota con potencia positiva descarta una excepción prefijada.}


\section{\B{An exact reflected sieve and batch certificate}{Criba reflejada exacta y certificado por lotes}}\label{sec:sieve}
\B{Put $C_a=3+2a$ and $s=N/2-3$. A distinct pair of odd integers at this target is uniquely $(C_a,C_{s-a})$ with $a\in I_N=\{a\in\Z:0\le a<s-a\}$. The midpoint is $m=N/2$ and its offset is $d=s-2a$, so the endpoints are $m-d,m+d$. Here $3$ is the origin of a progression of difference 2; it is neither a difference-3 progression, a ternary notation, nor a modulus. A wheel such as 210 is a separate operation.}{Sean $C_a=3+2a$ y $s=N/2-3$. Cada pareja de impares distintos para este objetivo se escribe de modo único como $(C_a,C_{s-a})$, con $a\in I_N=\{a\in\Z:0\le a<s-a\}$. El punto medio es $m=N/2$ y su desplazamiento es $d=s-2a$, de modo que los extremos son $m-d,m+d$. Aquí $3$ es el origen de una progresión de diferencia 2; no es una progresión de diferencia 3, una notación ternaria ni un módulo. Una rueda como la de 210 es una operación distinta.}

\B{Square-start criterion. For each odd prime $\ell$, set $u_\ell=(\ell^2-3)/2$ and $B_\ell=\{u_\ell+j\ell:j\in\Z_{\ge0}\}$. With $z=\lfloor\sqrt{N-3}\rfloor$, one has the exact identity}{Criterio de arranque en el cuadrado. Para cada primo impar $\ell$, sean $u_\ell=(\ell^2-3)/2$ y $B_\ell=\{u_\ell+j\ell:j\in\Z_{\ge0}\}$. Con $z=\lfloor\sqrt{N-3}\rfloor$, se tiene la identidad exacta}
\begin{equation}\label{eq:mark}
 R(N)=\#\left(I_N\setminus\bigcup_{\substack{3\le\ell\le z\\\ell\ \B{\mathrm{prime}}{\mathrm{primo}}}}\left(B_\ell\cup(s-B_\ell)\right)\right).
\end{equation}
\B{Indeed, $C_{u_\ell+j\ell}=\ell^2+2j\ell$ is composite. Conversely, an odd composite endpoint $n\le N-3$ has a least prime factor $\ell\le\sqrt n$, and $n/\ell$ is odd and at least $\ell$. Hence $n/\ell=\ell+2j$, so its index belongs to $B_\ell$. Apply the same proof to $s-a$ for the second endpoint; intersection with $I_N$ clips the two progressions in opposite directions. The empty interval at $N=4,6$ gives $R(N)=0$, not a counterexample to Goldbach.}{En efecto, $C_{u_\ell+j\ell}=\ell^2+2j\ell$ es compuesto. Recíprocamente, un extremo impar compuesto $n\le N-3$ tiene un menor factor primo $\ell\le\sqrt n$, y $n/\ell$ es impar y no menor que $\ell$. Luego $n/\ell=\ell+2j$ y su índice pertenece a $B_\ell$. Se aplica la misma prueba a $s-a$ para el segundo extremo; la intersección con $I_N$ recorta ambas progresiones en sentidos opuestos. El intervalo vacío para $N=4,6$ da $R(N)=0$, no un contraejemplo a Goldbach.}

\B{Batch criterion. Let $Q$ be the complete set of odd primes up to $z$ and $P=\prod_{\ell\in Q}\ell$ (empty product 1). For $n$ an odd integer with $3\le n\le N-3$, the exact test is}{Criterio por lotes. Sea $Q$ el conjunto completo de primos impares hasta $z$ y sea $P=\prod_{\ell\in Q}\ell$ (producto vacío igual a 1). Para $n$ impar con $3\le n\le N-3$, el criterio exacto es}
\[
 n\ \B{\text{prime}}{\text{primo}}\quad\Longleftrightarrow\quad n\in Q\ \text{or}\ \gcd(n,P)=1.
\]
\B{A composite cannot survive by the least-factor argument; a prime outside the batch is coprime to $P$. On the subset where neither endpoint belongs to $Q$, both endpoint tests combine into $\gcd(C_a C_{s-a},P)=1$. Each endpoint may be reduced modulo $P$ before multiplying. For the omitted small-endpoint cases first require that the corresponding endpoint equal a member of $Q$, then apply the individual test to the other. Omitting this exception can remove genuine pairs; completeness of $Q$ is essential to the converse. These are two proofs of the specification of the finite calculation, not an implication that the final survivor set is always nonempty.}{Un compuesto no puede sobrevivir por el argumento del menor factor; un primo ajeno al lote es coprimo con $P$. Cuando ninguno de los extremos pertenece a $Q$, las dos pruebas se agrupan en $\gcd(C_a C_{s-a},P)=1$. Cada extremo puede reducirse módulo $P$ antes de multiplicar. Para los casos de extremo pequeño, se exige primero que ese extremo sea igual a un miembro de $Q$ y después se aplica el criterio individual al otro. Omitir esta excepción puede eliminar parejas auténticas; la completitud de $Q$ es esencial para la implicación recíproca. Son dos pruebas de la especificación del cálculo finito, no una implicación de que siempre sobreviva alguna pareja.}

\B{Local factors. For odd $\ell$ the two forbidden index classes are $-3\,2^{-1}$ and $s+3\,2^{-1}$ modulo $\ell$. They coincide exactly when $\ell\mid N$. In endpoint coordinates, the classes are simply $0,N\pmod\ell$. At 2, for even $N$, there is one forbidden class. The Chinese remainder theorem therefore gives, with $\nu_\ell(N)=|\{0,N\}\bmod\ell|$,}{Factores locales. Para $\ell$ impar las dos clases prohibidas de índices son $-3\,2^{-1}$ y $s+3\,2^{-1}$ módulo $\ell$. Coinciden exactamente cuando $\ell\mid N$. En coordenadas de los extremos las clases son simplemente $0,N\pmod\ell$. En 2, para $N$ par, hay una clase prohibida. El teorema chino del resto da, con $\nu_\ell(N)=|\{0,N\}\bmod\ell|$,}
\[
 A_{210}(N)=\prod_{\ell\mid210}(\ell-\nu_\ell(N)),\qquad
 \Sg_{210}(N)=\frac{210 A_{210}(N)}{48^2}.
\]
\B{The first quantity counts classes over a full period, not pairs in $I_N$. It lies between 15 and 48. The finite wheel ignores larger factors and omits valid pairs containing 3, 5 or 7. Neither its positivity nor the positivity of an Euler product settles the finite-interval existence question.}{La primera cantidad cuenta clases en un periodo completo, no parejas en $I_N$. Está entre 15 y 48. La rueda finita ignora factores mayores y omite parejas válidas que contienen 3, 5 o 7. Ni su positividad ni la de un producto de Euler resuelven la existencia en el intervalo finito.}


\paragraph{\B{Algorithmic scope.}{Alcance algorítmico.}}
\B{The marker update is an Eratosthenes-style operation in the sense discussed by O'Neill \cite{oneill}. Through a bound $B$, the number of marking events is $O(B\log\log B)$, by summing $O(B/\ell)$ over odd primes $\ell\le\sqrt B$. The bound counts events, not the cost of the dictionary or priority queue, nor the work of scanning reflected targets. In a full period the wheel retains $48/210$ of the integer residue classes for a single operand; the $162/210$ discarded fraction is of classes, not of all composites. No speedup for counting all pairs is claimed.}{La actualización de marcadores es una operación de tipo Eratóstenes en el sentido discutido por O'Neill \cite{oneill}. Hasta una cota $B$, hay $O(B\log\log B)$ eventos de marcado, sumando $O(B/\ell)$ sobre los primos impares $\ell\le\sqrt B$. Se cuentan eventos, no el coste del diccionario o de la cola de prioridad ni el recorrido de los objetivos reflejados. En un periodo completo, la rueda retiene $48/210$ de las clases residuales enteras de un solo operando; la fracción descartada $162/210$ corresponde a clases, no a todos los compuestos. No se afirma una aceleración del recuento de todas las parejas.}


\section{\B{New calculation on the published 1000-prime prefix}{Cálculo nuevo sobre el prefijo publicado de mil primos}}
\B{We used every even $N$ from 4 through 7930 inclusive. For distinct odd endpoints the larger endpoint is at most $N-3\le7927$. The additive generator returned the same 1000 values as the table in \cite[Appendix B]{additive}, checked entry by entry against the PDF text. The run compared the sets of smaller endpoints produced by square-start reflection, complete-batch GCD with the exceptions, and direct reflection of an independently constructed Eratosthenes table. No probabilistic primality tests or floating-point operations enter this comparison. This is a transfer experiment, not a new computational bound for Goldbach.}{Se usaron todos los pares $N$ desde 4 hasta 7930 inclusive. Para extremos impares distintos, el mayor no excede $N-3\le7927$. El generador aditivo devolvió los mismos mil valores de la tabla de \cite[apéndice B]{additive}, comprobados entrada a entrada frente al texto del PDF. Se compararon los conjuntos del extremo menor obtenidos por reflexión con arranque en el cuadrado, MCD por lote completo con las excepciones y reflexión directa sobre una tabla independiente de Eratóstenes. En la comparación no intervienen tests probabilísticos de primalidad ni operaciones en coma flotante. Es un experimento de transferencia, no una nueva cota computacional para Goldbach.}

\begin{center}\small
\begin{tabular}{@{}lrrrrr@{}}
\toprule
\B{Domain}{Dominio}&\B{Targets}{Objetivos}&$\sum R(N)$&$\min R(N)$&\B{False positives}{Falsos positivos}&\B{Omitted pairs}{Parejas omitidas}\\
\midrule
\B{All even targets}{Todos los pares}&3964&283488&0&538737&2992\\
$H$&2503&186315&2&351829&1944\\
$H_1$&913&72916&1&139347&620\\
\bottomrule
\end{tabular}
\end{center}
\B{The last two columns refer to screening by $\gcd(pq,210)=1$ alone; totals are target--pair incidences. The two zero targets for the distinct-pair count are 4 and 6, both covered by a prime diagonal. The wheel rejects both $(3,13)$ and $(5,11)$ at $16\in H$. It admits $(11,121)$ at 132; the complete divisor batch rejects it. All three exact algorithms agree on every target, but no inference beyond the specified interval is drawn. Source code, full aggregate output and verification instructions accompany this manuscript.}{Las dos últimas columnas se refieren al filtrado exclusivo por $\gcd(pq,210)=1$; los totales cuentan incidencias objetivo--pareja. Los dos objetivos con recuento de parejas distintas cero son 4 y 6, ambos cubiertos por una diagonal prima. La rueda rechaza tanto $(3,13)$ como $(5,11)$ en $16\in H$. Admite $(11,121)$ en 132; el lote completo de divisores la descarta. Los tres algoritmos exactos concuerdan en cada objetivo, pero no se infiere nada fuera del intervalo declarado. El código, la salida agregada completa y las instrucciones de verificación acompañan a este manuscrito.}


\paragraph{\B{Larger witness-only audit.}{Comprobación de testigos en un intervalo mayor.}}
\B{A separate C++17 program checks every even target through $8\cdot10^6$ by exact Eratosthenes and linear least-prime-factor sieves. It compares the primality tables and checks agreement of the least strictly smaller endpoint found for each target, using two separate scans. This is a coverage check, not a count of all representations.}{Un programa separado en C++17 comprueba todos los pares hasta $8\cdot10^6$ mediante Eratóstenes y una criba lineal de menor factor primo, ambas exactas. Compara las tablas de primalidad y contrasta el menor extremo estrictamente inferior encontrado para cada objetivo mediante dos recorridos distintos. Es una prueba finita de cobertura, no un recuento de todas las representaciones.}
\begin{center}\small
\begin{tabular}{@{}lrr@{}}
\toprule
\B{Domain}{Dominio}&\B{Targets}{Objetivos}&\B{With a distinct-prime witness}{Con testigo de primos distintos}\\
\midrule
$H$&3\,202\,369&3\,202\,369\\
$H_1$&514\,484&514\,484\\
\B{Prime midpoint}{Punto medio primo}&283\,146&283\,144\\
\bottomrule
\end{tabular}
\end{center}
\B{Only $4$ and $6$ lack a distinct pair in this interval; their prime diagonals are $2+2$ and $3+3$. This does not make the three cases disjoint.}{Solo $4$ y $6$ carecen de pareja distinta en el intervalo; sus diagonales primas son $2+2$ y $3+3$. Esto no convierte los tres casos en disjuntos.}


\section{\B{Why local density and global energy do not suffice}{Por qué la densidad local y la energía global no bastan}}
\B{The project considered successive routes from averaged or local data to a prescribed target. We record the mathematical reasons for stopping them. These obstructions are statements about specific methods or information classes, not counterexamples to Goldbach.}{El proyecto examinó distintas vías para pasar de datos medios o locales a un objetivo prefijado. Se recogen los motivos matemáticos para detenerlas. Las obstrucciones conciernen a métodos o clases de información concretos, no son contraejemplos a Goldbach.}

\paragraph{\B{Variance and correlations (stages 56--63).}{Varianza y correlaciones (etapas 56--63).}}
\B{For a block of $m\ge2$ nonnegative values with mean $\bar v$, one zero forces population variance at least $\bar v^2/(m-1)$: Cauchy--Schwarz bounds the sum of the other $m-1$ squares below by $m^2\bar v^2/(m-1)$. The strict reverse inequality is a sufficient certificate, not a necessary condition for positivity. For example, $(1,1,10)$ fails the strict test despite containing no zero. Positive shifted products on average do not rule out isolated zeros, as the sequence $(0,1,1)$ already shows. For cyclic energy $E=\sum_j(v_{j+1}-v_j)^2$, Fourier diagonalization and Cauchy--Schwarz give the sharp bound $|v_i-\bar v|^2\le(m^2-1)E/(12m)$ for $m\ge3$; the constant is the mean of the nonzero inverse eigenvalues $[4\sin^2(\pi j/m)]^{-1}$, and equality is attained on a column of the Laplacian pseudoinverse. For the three-point path the squared-mean minus $2E/9$ reads $(-a^2-3b^2-c^2+6ab+2ac+6bc)/9$; if positive with nonnegative entries then $b>0$, since at $b=0$ the expression is $-(a-c)^2/9$. These tests need actual arithmetic control of the inputs.}{Para un bloque de $m\ge2$ valores no negativos de media $\bar v$, un cero obliga a que la varianza poblacional sea al menos $\bar v^2/(m-1)$: Cauchy--Schwarz acota por abajo la suma de los otros $m-1$ cuadrados por $m^2\bar v^2/(m-1)$. La desigualdad estricta inversa es un criterio suficiente, no necesario, de positividad. Por ejemplo, $(1,1,10)$ no pasa el criterio estricto pese a no contener ceros. Los productos desplazados positivos en promedio no excluyen ceros aislados, como muestra $(0,1,1)$. Para la energía cíclica $E=\sum_j(v_{j+1}-v_j)^2$, la diagonalización de Fourier y Cauchy--Schwarz dan la cota óptima $|v_i-\bar v|^2\le(m^2-1)E/(12m)$ para $m\ge3$; la constante es la media de los inversos de autovalores no nulos $[4\sin^2(\pi j/m)]^{-1}$ y se alcanza sobre una columna de la pseudoinversa del laplaciano. En el camino de tres puntos, el cuadrado de la media menos $2E/9$ es $(-a^2-3b^2-c^2+6ab+2ac+6bc)/9$; si es positivo con entradas no negativas, entonces $b>0$, pues para $b=0$ la expresión es $-(a-c)^2/9$. Estos criterios requieren control aritmético de sus entradas.}

\B{The products $R(n)R(m)$ split into pairs of representations with one common prime or with four distinct primes for $n\ne m$, since two common primes would force equal targets. The shared-prime term is the count of $x$ with $x,n-x,m-x$ prime and both diagonals excluded. For $(12,18)$ the shared positions are $5,7$; $11$ is not one because $12-11=1$ (the stage 64 erratum). At the same target the identity is $R(n)^2=R(n)+R(n)(R(n)-1)$, with the off-diagonal term having four distinct prime values. These decompositions do not supply lower bounds for the correlations.}{Los productos $R(n)R(m)$ se descomponen en pares de representaciones con un primo común o cuatro primos distintos si $n\ne m$, ya que dos primos comunes forzarían objetivos iguales. El término de primo común cuenta los $x$ para los que $x,n-x,m-x$ son primos, excluidas ambas diagonales. Para $(12,18)$ los valores comunes son $5,7$; no lo es $11$ porque $12-11=1$ (corrección de la etapa 64). Para el mismo objetivo se tiene $R(n)^2=R(n)+R(n)(R(n)-1)$, con cuatro valores primos distintos en el término no diagonal. Las descomposiciones no aportan cotas inferiores para las correlaciones.}

\B{Even at the level of the arithmetic singular series $\Sg$ defined below, wheel normalization does not imply smoothness. For a finite set of primes $\ell\ge11$, Dirichlet's theorem and CRT provide arbitrarily large primes $q\equiv1\pmod\ell$ for all these primes and $q\equiv4\pmod5$. Then $N=2(q-1)\in H$ since $q-1$ is even composite and $2q-3$ is a multiple of 5 greater than 5. The selected primes divide $N$, whereas $N+2=2q$. The factors $(\ell-1)/(\ell-2)$ over growing finite sets have unbounded product by divergence of the sum of reciprocal primes. Since $15\le A_{210}\le48$, the ratio of the wheel-normalized singular series at these consecutive even arguments is unbounded. This is not a claim about actual prime-pair asymptotics.}{Incluso para la serie singular aritmética $\Sg$ definida abajo, la normalización por la rueda no implica suavidad. Dado un conjunto finito de primos $\ell\ge11$, el teorema de Dirichlet y el teorema chino del resto proporcionan primos arbitrariamente grandes $q\equiv1\pmod\ell$ para todos ellos y $q\equiv4\pmod5$. Entonces $N=2(q-1)\in H$ porque $q-1$ es par compuesto y $2q-3$ es un múltiplo de 5 mayor que 5. Los primos elegidos dividen $N$, mientras que $N+2=2q$. Los factores $(\ell-1)/(\ell-2)$ sobre conjuntos finitos crecientes tienen producto no acotado por divergencia de la suma de inversos de los primos. Como $15\le A_{210}\le48$, el cociente de las series normalizadas en esos pares consecutivos no está acotado. Esto no afirma una asintótica de parejas primas reales.}

\paragraph{\B{Fourier localization (stages 64--70).}{Localización de Fourier (etapas 64--70).}}
\B{For $S(\alpha)=\sum_p a_p\e(p\alpha)$, where $\e(t)=\exp(2\pi it)$, the $N$th coefficient of $S^2$ counts sums; $|S|^2$ counts differences. If $c(n)$ is the $n$th coefficient of $\one_E S^2$ on a measurable set of arcs $E$, Parseval gives $\sum_n|c(n)|^2=\int_E |S|^4$. A budget $B$ allows as many as $B/A^2$ coefficients of size at least $A$; if $B\ge A^2$, a singleton of exactly this size is already admissible in the abstract coefficient class. An assumed increment bound $|c(N+h)-c(N)|\le d_h$ for distinct nonzero $h$ raises the necessary energy to $A^2+\sum_h\pos{A-d_h}^2$. The triangle inequality proves the bound, and putting all of the selected entries on the same ray and setting all others to zero attains it. With $m$ shifts and a squared-increment budget $V$, the corresponding bound is $A^2+m\pos{A-\sqrt{V/m}}^2$, by the reverse triangle inequality in Euclidean norm. Equal increments show sharpness for the upper-budget problem. These are deterministic statements, not evidence that prime coefficients realize the extremizers.}{Para $S(\alpha)=\sum_p a_p\e(p\alpha)$, con $\e(t)=\exp(2\pi it)$, el coeficiente $N$ de $S^2$ cuenta sumas; $|S|^2$ cuenta diferencias. Si $c(n)$ es el coeficiente $n$ de $\one_E S^2$ para un conjunto medible de arcos $E$, Parseval da $\sum_n|c(n)|^2=\int_E|S|^4$. Un presupuesto $B$ admite hasta $B/A^2$ coeficientes de módulo al menos $A$; si $B\ge A^2$, ya admite un coeficiente aislado de ese tamaño en la clase abstracta. Una cota supuesta de incrementos $|c(N+h)-c(N)|\le d_h$ para distintos $h\ne0$ aumenta la energía necesaria a $A^2+\sum_h\pos{A-d_h}^2$. La desigualdad triangular prueba la cota; situar las entradas elegidas en el mismo rayo y anular las demás alcanza la igualdad. Con $m$ desplazamientos y presupuesto $V$ de incrementos al cuadrado, se obtiene $A^2+m\pos{A-\sqrt{V/m}}^2$ por la desigualdad triangular inversa en norma euclídea. Incrementos iguales muestran optimalidad para el problema de presupuesto superior. Son enunciados deterministas, no pruebas de que los coeficientes de los primos realicen los extremizadores.}

\B{Arc restriction retains more geometry. For shifts in a fixed finite set, let $G_E$ be the Gram matrix of $\e(-h\alpha)-1$ on $E$. The target factor has modulus one, so it cancels from the Gram matrix. For disjoint pieces $E_j$ and independent complex $L^2$ budgets $B_j$, the largest possible norm of an increment vector over that unrestricted class is $\max_{\|z\|=1}\sum_j\sqrt{B_j}\sqrt{z^*G_{E_j}z}$: take the dual norm and choose each input parallel to the corresponding adjoint image. With just a total budget, maximizing over its allocation by Cauchy--Schwarz recovers the single Gram bound on the union. A rational center $a/q$ has zero shift multiplier only when $q\mid h$ in lowest terms. Neither this geometric structure nor the project's full-circle nonnegative-coefficient cone relaxation yields an actual-prime signed estimate. The latter relaxation drops the nonlinear condition that the polynomial is the square of the full prime sum.}{La restricción a arcos retiene más geometría. Para un conjunto fijo finito de desplazamientos, sea $G_E$ la matriz de Gram de $\e(-h\alpha)-1$ en $E$. El factor del objetivo tiene módulo uno y desaparece de la matriz. Para piezas disjuntas $E_j$ y presupuestos complejos $L^2$ independientes $B_j$, la máxima norma de un vector de incrementos en esa clase sin restricciones es $\max_{\|z\|=1}\sum_j\sqrt{B_j}\sqrt{z^*G_{E_j}z}$: se toma la norma dual y se elige cada entrada paralela a la imagen adjunta correspondiente. Con un único presupuesto total, maximizar su reparto por Cauchy--Schwarz devuelve la cota de Gram sobre la unión. En un centro racional irreducible $a/q$ el multiplicador del desplazamiento solo se anula si $q\mid h$. Ni esta geometría ni la relajación del proyecto al cono de coeficientes no negativos en todo el círculo dan una estimación con signo para los primos auténticos. La segunda relajación omite la condición no lineal de que el polinomio sea el cuadrado de la suma completa de primos.}

\paragraph{\B{Deletion and positive arcs (stages 70--73).}{Borrado y arcos positivos (etapas 70--73).}}
\B{For $N\in D$, complementary prime pairs form a matching because $p$ has at most one partner $N-p$ and there is no diagonal. With nonnegative weights, deleting the lighter endpoint of each active edge annihilates that target coefficient in the modified support. The exact minimum deleted mass is $\sum_{\{p,q\}}\min(a_p,a_q)$: one endpoint per edge is necessary and different edges are disjoint. This is not valid on arbitrary graphs: a two-edge star with unit weights is erased by deleting its center at cost 1, not 2. Using a standard upper-bound sieve \cite{ford,wu} and $N/\varphi(N)\ll\log\log N$ gives $R(N)\ll N\log\log N/(\log N)^2$ for large even $N$. For the smooth prime weights below and $N\in[X/2,X]\cap D$, the removed mass $\delta$ is therefore $O_w(X\log\log X/\log X)=o(X)$. The retained first-order total and fixed reduced-class asymptotics follow from the prime number theorem in fixed progressions. Also, nonnegative sum coefficients can only decrease, so the full fourth moment of the modified sum cannot increase. These changes do not preserve the exact set of primes.}{Para $N\in D$, las parejas primas complementarias son aristas disjuntas porque $p$ tiene a lo sumo un compañero $N-p$ y no hay diagonal. Con pesos no negativos, borrar el extremo más ligero de cada arista activa anula el coeficiente de ese objetivo en el soporte modificado. La masa mínima borrada es exactamente $\sum_{\{p,q\}}\min(a_p,a_q)$: es necesario quitar un extremo por arista y las aristas distintas no se tocan. No vale para grafos arbitrarios: una estrella de dos aristas con pesos unidad se anula borrando su centro a coste 1, no 2. Una criba superior estándar \cite{ford,wu} y $N/\varphi(N)\ll\log\log N$ dan $R(N)\ll N\log\log N/(\log N)^2$ para pares grandes. Para los pesos suaves de abajo y $N\in[X/2,X]\cap D$, la masa borrada $\delta$ es, por tanto, $O_w(X\log\log X/\log X)=o(X)$. El teorema de los números primos en progresiones fijas conserva las asintóticas de primer orden de la masa y de cada clase reducida fija. Además, los coeficientes no negativos de las sumas solo pueden disminuir, de modo que el cuarto momento completo del polinomio modificado no crece. Estos cambios no preservan el conjunto exacto de los primos.}

\B{In particular, for nonnegative $a_p$ supported in $(0,X)$ write $M=\sum a_p$. At integers $N\in[X/2,X]$, integration over $[-1/(6X),1/(6X)]$ contributes at least $2M^2/(3\pi X)$ to the real sum coefficient: each term has $|p+q-N|\le3X/2$, so the integrated sinc is at least $2/(3\pi X)$. The prime number theorem makes $M\asymp_w X$ for the full smooth sum and deletion leaves its leading term unchanged. More generally on a measurable arc set of length $|E|$, changing the weights by downward attenuation of mass $\delta$ alters its target integral by at most $|E|\delta(2M-\delta)$, by factoring the difference of squares. Thus fixed cores of total length $O(1/X)$ retain their leading profiles under this deletion while the full coefficient can be zero. This rules out a positivity proof using only the retained data across the enlarged subset class. It is neither a parity-breaking estimate nor a claim that the exact Vaughan identity or maximal Bombieri--Vinogradov precision holds for the modified sequence.}{En particular, para $a_p\ge0$ con soporte en $(0,X)$ sea $M=\sum a_p$. En objetivos enteros $N\in[X/2,X]$, la integral en $[-1/(6X),1/(6X)]$ aporta al coeficiente real al menos $2M^2/(3\pi X)$: cada término tiene $|p+q-N|\le3X/2$, luego el seno cardinal integrado es al menos $2/(3\pi X)$. El teorema de los números primos da $M\asymp_w X$ para la suma suave completa y el borrado no altera su término principal. Más en general, en un conjunto de arcos de longitud $|E|$, atenuar hacia abajo los pesos con masa $\delta$ modifica la integral del objetivo en a lo sumo $|E|\delta(2M-\delta)$, al factorizar la diferencia de cuadrados. Así, los núcleos fijos de longitud total $O(1/X)$ conservan sus perfiles principales tras el borrado, mientras que el coeficiente completo puede valer cero. Se descarta una prueba de positividad basada solo en los datos conservados para toda la clase ampliada de subconjuntos. No es una estimación que rompa la barrera de paridad, ni se afirma que la identidad exacta de Vaughan o la precisión maximal de Bombieri--Vinogradov valgan para la sucesión modificada.}

\section{\B{The Type I reduction and its kernel}{La reducción de Tipo I y su núcleo}}
\B{T73.5 and T74.3 below are inherited analytical derivations using classical inputs, presented with proofs for scrutiny, not independently accepted theorems of the project. T75.1 is the error transport that depends on them.}{T73.5 y T74.3 son derivaciones analíticas heredadas con ingredientes clásicos, expuestas con pruebas para su examen, no teoremas del proyecto aceptados de forma independiente. T75.1 es el transporte de errores que depende de ellas.}

\B{Fix a nonnegative $w\in C_c^\infty((0,1))$ with support in a fixed $[\sigma,\rho]\subset(0,1)$. Put $h_N(t)=w(t/X)w((N-t)/X)$ and set $U=V=\max(1,\lfloor X^{1/5}\rfloor)$. The actual prime sum and the ordered weighted coefficients are}{Fijamos $w\ge0$ en $C_c^\infty((0,1))$ con soporte en un intervalo fijo $[\sigma,\rho]\subset(0,1)$. Sea $h_N(t)=w(t/X)w((N-t)/X)$ y sean $U=V=\max(1,\lfloor X^{1/5}\rfloor)$. La suma sobre los primos auténticos y los coeficientes ponderados ordenados son}
\begin{align*}
 S_X(\alpha)&=\sum_{p\ \B{\mathrm{prime}}{\mathrm{primo}}}(\log p)w(p/X)\e(p\alpha),\qquad r_X(N)=[S_X^2]_N,\\
 G_X(N)&=\sum_{n\ge1}\Lambda(n)\Lambda(N-n)h_N(n).
\end{align*}
\B{Extend $\Lambda$ by zero on nonpositive integers. Write $b_U(k)=\sum_{d\mid k,\ d>U}\mu(d)$ and let $c_l=\sum_{dm=l,\ d\le U,\ m\le V}\mu(d)\Lambda(m)$, define}{Extendemos $\Lambda$ por cero en los enteros no positivos. Para $b_U(k)=\sum_{d\mid k,\ d>U}\mu(d)$ y $c_l=\sum_{dm=l,\ d\le U,\ m\le V}\mu(d)\Lambda(m)$, se define}
\begin{align*}
 A_3(N)&=\sum_{\substack{k>U,\ m>V\\km<N}} b_U(k)\Lambda(m)\Lambda(N-km)h_N(km),\\
 \mathcal K(t)&=\sum_{d\le U,\ (d,N)=1}\frac{\mu(d)}{\varphi(d)}\log(t/d)
       -\sum_{l\le UV,\ (l,N)=1}\frac{c_l}{\varphi(l)},\qquad
 \mathcal M_I(N)=\int h_N(t)\mathcal K(t)\,dt .
\end{align*}

\paragraph{T73.5.}
\B{For every fixed $A>0$, uniformly in integers $N\in[X/2,X]$ for sufficiently large $X$,}{Para cada $A>0$ fijo, uniformemente para enteros $N\in[X/2,X]$ y $X$ suficientemente grande,}
\begin{equation}\label{eq:TI}
 G_X(N)=\mathcal M_I(N)+A_3(N)+E_I(N),\qquad
 E_I(N)\ll_{A,w}X\log^{-A}X+X^{2/5}\log^3 X.
\end{equation}
\B{Proof. With Dirichlet convolution $*$ and subscripts denoting upper or lower truncations, the exact identity is}{Prueba. Con la convolución de Dirichlet $*$ y subíndices para truncamientos superiores o inferiores, la identidad exacta es}
\[
 \Lambda=\Lambda_{\le V}+\mu_{\le U}*\log
         -\mu_{\le U}*\Lambda_{\le V}*\one+\mu_{>U}*\one*\Lambda_{>V}.
\]
\B{Here $\one(n)=1$. It follows from $\log=\one*\Lambda$ and $\mu*\one=\delta_1$, with no assumptions on pairs. In $G_X$ the term $\Lambda_{\le V}$ vanishes for $V<\sigma X$. The last term is $A_3$, since $\mu_{>U}*\one$ vanishes for $k\le U$. In the remaining terms, $n=dk$ or $n=lk$, respectively. Change variables to $x=N-n$; the weights are $h_N(N-x)\log((N-x)/d)$ and $h_N(N-x)$, with total variations $O_w(\log X)$ and $O_w(1)$, and support away from the endpoints. Also $|c_l|\le\sum_{m\mid l}\Lambda(m)=\log l$. On a reduced class $x\equiv N\pmod d$ or $\pmod l$, Stieltjes partial summation against $d\psi(x;l,N)-dx/\varphi(l)$ costs the maximal discrepancy times the variation. The maximal classical Bombieri--Vinogradov estimate \cite{mv,bv} at $Q=UV\le X^{2/5}$ absorbs the extra logarithm by choosing its saving exponent larger than $A$. Its standard main term is $y/\varphi(l)$; no unproved zero-free region is required. On a nonreduced class, a contributing $x$ must be a power of a prime dividing the modulus and $N$; each such prime base contributes at most $\log X$ of $\Lambda$-mass. Thus the nonreduced contribution is $O_w(UV\log^3 X)$ after including the other log weight. The reduced main terms are exactly $\mathcal M_I$.}{Aquí $\one(n)=1$. Se deduce de $\log=\one*\Lambda$ y $\mu*\one=\delta_1$ sin hipótesis sobre parejas. En $G_X$ el término $\Lambda_{\le V}$ se anula para $V<\sigma X$. El último término da $A_3$, porque $\mu_{>U}*\one$ se anula para $k\le U$. En los restantes términos, $n=dk$ o $n=lk$, respectivamente. Cambiamos a $x=N-n$; los pesos son $h_N(N-x)\log((N-x)/d)$ y $h_N(N-x)$, con variaciones totales $O_w(\log X)$ y $O_w(1)$ y soporte alejado de los extremos. Además, $|c_l|\le\sum_{m\mid l}\Lambda(m)=\log l$. En una clase reducida $x\equiv N\pmod d$ o $\pmod l$, la sumación parcial de Stieltjes frente a $d\psi(x;l,N)-dx/\varphi(l)$ cuesta la discrepancia maximal por la variación. La versión maximal clásica de Bombieri--Vinogradov \cite{mv,bv} con $Q=UV\le X^{2/5}$ absorbe el logaritmo adicional al elegir su exponente de ahorro mayor que $A$. Su término principal estándar es $y/\varphi(l)$; no se necesita una región libre de ceros no demostrada. En una clase no reducida, un $x$ que contribuya es potencia de un primo divisor del módulo y de $N$; cada base prima aporta a lo sumo $\log X$ de masa de $\Lambda$. De ahí la contribución no reducida $O_w(UV\log^3 X)$, incluido el otro peso logarítmico. Los términos principales reducidos son exactamente $\mathcal M_I$.}

\paragraph{T74.3.}
\B{For positive even $N\le X$ and $t\in[\sigma X,\rho X]$, uniformly as $X\to\infty$,}{Para $N$ par positivo, $N\le X$ y $t\in[\sigma X,\rho X]$, uniformemente cuando $X\to\infty$,}
\begin{align}
 \mathcal K(t)&=\Sg(N)+O_{\sigma,\rho}(\e^{-c\sqrt{\log X}}),\qquad c>0,\label{eq:kernel}\\
 \Sg(N)&=2C_2\!\prod_{p\mid N,\ p>2}\frac{p-1}{p-2},\quad
 C_2=\prod_{p>2}\left(1-\frac{1}{(p-1)^2}\right)>0.\nonumber
\end{align}
\B{Consequently $\mathcal M_I=F_X+E_K$, where $F_X(N)=X\Sg(N)J_w(N/X)$, $J_w(\tau)=\int w(u)w(\tau-u)\,du$ and $E_K=O_w(X\e^{-c\sqrt{\log X}})$. This integrates the kernel uniformly over support of length $O(X)$.}{Por tanto $\mathcal M_I=F_X+E_K$, donde $F_X(N)=X\Sg(N)J_w(N/X)$, $J_w(\tau)=\int w(u)w(\tau-u)\,du$ y $E_K=O_w(X\e^{-c\sqrt{\log X}})$. Se integra el núcleo uniformemente sobre un soporte de longitud $O(X)$.}

\B{Proof of the kernel evaluation. Define $a_N(n)=\mu(n)\one_{(n,N)=1}/\varphi(n)$ and $b(n)=\mu(n)/n$. Euler factors give $a_N=b*g_N$, where $g_N$ is multiplicative and for $j\ge1$}{Prueba de la evaluación del núcleo. Definimos $a_N(n)=\mu(n)\one_{(n,N)=1}/\varphi(n)$ y $b(n)=\mu(n)/n$. Los factores de Euler dan $a_N=b*g_N$, donde $g_N$ es multiplicativa y para $j\ge1$}
\[
 g_N(p^j)=\begin{cases}p^{-j},&p\mid N,\\-p^{-j}/(p-1),&p\nmid N.\end{cases}
\]
\B{For $N\le X^B$ with $B$ fixed, $T_N=\sum_k|g_N(k)|k^{1/4}\le\exp(O_B((\log X)^{1/4}))$. The factors for $p\nmid N$ are uniformly summable in their logarithms. For the others, set $L=\log X$: the sum of $p^{-3/4}$ at $p\le L$ is $O(L^{1/4})$, while at $p>L$ at most $BL/\log L$ primes can divide $N$, since their product is at most $X^B$. This bound also holds with any one local factor removed.}{Para $N\le X^B$ con $B$ fijo, $T_N=\sum_k|g_N(k)|k^{1/4}\le\exp(O_B((\log X)^{1/4}))$. Los logaritmos de los factores para $p\nmid N$ son uniformemente sumables. Para el resto, sea $L=\log X$: la suma de $p^{-3/4}$ para $p\le L$ es $O(L^{1/4})$, mientras que para $p>L$ pueden dividir a $N$ a lo sumo $BL/\log L$ primos, pues su producto no excede $X^B$. La misma cota vale si se quita un factor local.}

\B{From the classical quantitative Mertens estimate \cite[Ch.~6]{mv}, partial summation and the expansion $1/\zeta(1+s)=s+O(s^2)$ at zero,}{Por la estimación cuantitativa clásica para la función sumatoria de Möbius \cite[cap.~6]{mv}, sumación parcial y el desarrollo $1/\zeta(1+s)=s+O(s^2)$ en cero,}
\[
 \sum_{n\le y} b(n)=O(\e^{-c_0\sqrt{\log y}}),\qquad
 \sum_{n\le y} b(n)\log n=-1+O(\e^{-c_0\sqrt{\log y}}).
\]
\B{The summatory estimate $\sum_{n\le y}\mu(n)\ll y\e^{-c_1\sqrt{\log y}}$ gives convergent partial-summation tails also after multiplying by $\log n$, uniformly for real $s\ge0$ near zero; this justifies evaluation and differentiation at the boundary. Convolve the last two displayed sums with $g_N$, using $\log(km)=\log k+\log m$ in the second. For $k\le\sqrt{y}$ insert the decay bound at $y/k\ge\sqrt y$; for $k>\sqrt y$, the absolute tail is at most $y^{-1/8}T_N$. Logarithmic factors are absorbed by reducing the positive exponent constant. For $y$ in a fixed polynomial range $X^\alpha\le y\le X^\beta$, $0<\alpha\le\beta$, this gives uniformly}{La cota $\sum_{n\le y}\mu(n)\ll y\e^{-c_1\sqrt{\log y}}$ da colas convergentes por sumación parcial también tras multiplicar por $\log n$, uniformemente para $s$ real no negativo cerca de cero; queda justificada la evaluación y derivación en la frontera. Convolucionamos las dos sumas mostradas con $g_N$, usando $\log(km)=\log k+\log m$ en la segunda. Para $k\le\sqrt{y}$ se aplica el decaimiento en $y/k\ge\sqrt y$; para $k>\sqrt y$, la cola absoluta no supera $y^{-1/8}T_N$. Los logaritmos se absorben reduciendo la constante positiva del exponente. Para $y$ en un rango polinómico fijo $X^\alpha\le y\le X^\beta$, $0<\alpha\le\beta$, se obtiene uniformemente}
\begin{align*}
 \sum_{n\le y}a_N(n)&=O(\e^{-c\sqrt{\log X}}),\\
 \sum_{n\le y}a_N(n)\log n&=-\sum_k g_N(k)+O(\e^{-c\sqrt{\log X}}).
\end{align*}
\B{The absolutely convergent product $\sum g_N(k)$ has factor $p/(p-1)$ at $p\mid N$ and $1-(p-1)^{-2}$ elsewhere, giving $\Sg(N)$ for even $N$. It remains to treat the second sum in $\mathcal K$. Each nonzero $\Lambda(m)$ has $m=p^j$ with $p\nmid N$. Let $g_{N,p}$ be $g_N$ with the $p$-factor removed, i.e., it vanishes on multiples of $p$. For squarefree $d$ divisible by $p$, $\varphi(dp^j)=p^j\varphi(d)$, not $\varphi(d)\varphi(p^j)$. The two squarefree exponents of $p$ in $d$ then give the identity}{El producto absolutamente convergente $\sum g_N(k)$ tiene factor $p/(p-1)$ para $p\mid N$ y $1-(p-1)^{-2}$ en el resto, lo que da $\Sg(N)$ para $N$ par. Queda la segunda suma de $\mathcal K$. Cada $\Lambda(m)$ no nulo tiene $m=p^j$ con $p\nmid N$. Sea $g_{N,p}$ la función $g_N$ sin el factor local en $p$, es decir, nula en múltiplos de $p$. Para $d$ libre de cuadrados divisible por $p$, $\varphi(dp^j)=p^j\varphi(d)$, no $\varphi(d)\varphi(p^j)$. Los dos exponentes posibles de $p$ en $d$ dan entonces}
\[
 \sum_{d\le U,\ (d,N)=1}\frac{\mu(d)}{\varphi(dp^j)}
 =\frac{1}{\varphi(p^j)}\sum_{d\le U}(b*g_{N,p})(d).
\]
\B{Indeed the normalized $p$-factor on the left is $1-p^{-1-s}$ in the Dirichlet series variable $s$, exactly that of $b$. The same moment and convolution argument bounds the last sum uniformly in $p$ by $O(\e^{-c\sqrt{\log X}})$. Finally $\sum_{p^j\le V}\log p/\varphi(p^j)\ll\log^2(2V)$ by summing the geometric tails and bounding by the corresponding integer sum. Therefore the second sum in the kernel tends uniformly to zero at the asserted scale, after adjusting $c$, as does $(\log t)\sum_{d\le U}a_N(d)$. The remaining term is $-\sum_{d\le U}a_N(d)\log d=\Sg(N)+O(\e^{-c\sqrt{\log X}})$. For a finite check of the shared-prime term, $N=16$, $U=V=3$ gives $\sum c_l/\varphi(l)=(1/2-1/6)\log3$, not $(1/2-1/4)\log3$. This is why naive totient separation fails.}{En efecto, el factor normalizado en $p$ de la izquierda es $1-p^{-1-s}$ en la variable de Dirichlet $s$, exactamente el de $b$. El mismo momento y argumento de convolución acota la última suma uniformemente en $p$ por $O(\e^{-c\sqrt{\log X}})$. Por último, $\sum_{p^j\le V}\log p/\varphi(p^j)\ll\log^2(2V)$ al sumar las colas geométricas y acotar mediante la suma sobre enteros. La segunda suma del núcleo tiende a cero uniformemente a la escala indicada, ajustando $c$, al igual que $(\log t)\sum_{d\le U}a_N(d)$. Queda $-\sum_{d\le U}a_N(d)\log d=\Sg(N)+O(\e^{-c\sqrt{\log X}})$. Como comprobación finita del caso de factor compartido, $N=16$, $U=V=3$ da $\sum c_l/\varphi(l)=(1/2-1/6)\log3$, no $(1/2-1/4)\log3$. Esa es la razón por la que falla la separación ingenua de la función de Euler.}

\section{\B{Error transport and the first missing theorem}{Transporte de errores y primer teorema que falta}}\label{sec:stop}
\paragraph{\B{Removing proper prime powers.}{Eliminación de las potencias propias de primos.}}
\B{Set $M_w=\|w\|_\infty$ and let $P_X(N)=G_X(N)-r_X(N)$. Nonnegativity and a union bound on which endpoint is a proper power give, for $X\ge4$,}{Sean $M_w=\|w\|_\infty$ y $P_X(N)=G_X(N)-r_X(N)$. La no negatividad y la cota por la unión de los casos en que un extremo es potencia propia dan, para $X\ge4$,}
\[
 0\le P_X(N)\le 2 M_w^2\pi(\sqrt X)(\log X)^2.
\]
\B{For each prime base $p\le\sqrt X$, the total proper-power $\Lambda$-mass is at most $\log X$; the reflected endpoint is uniquely determined and its weight is at most $M_w\log X$. Counting both orientations yields the factor 2. With $\pi(y)\le C_{\rm Ch}y/\log y$ for $y\ge4$, this gives $P_X\le 4C_{\rm Ch}M_w^2\sqrt X\log X$ uniformly in integer $N$ for $X\ge16$.}{Para cada base prima $p\le\sqrt X$, la masa total de $\Lambda$ de sus potencias propias no supera $\log X$; el extremo reflejado está determinado y su peso es a lo sumo $M_w\log X$. Contar ambas orientaciones da el factor 2. Con $\pi(y)\le C_{\rm Ch}y/\log y$ para $y\ge4$, se obtiene $P_X\le 4C_{\rm Ch}M_w^2\sqrt X\log X$ uniformemente para todo $N$ entero y $X\ge16$.}

\paragraph{T75.1.}
\B{Substituting the two Type I identities into $r_X=G_X-P_X$ gives}{Al sustituir las dos identidades de Tipo I en $r_X=G_X-P_X$ se obtiene}
\begin{equation}\label{eq:transport}
 A_3(N)=r_X(N)-F_X(N)+e_X(N),\quad
 e_X=P_X-E_K-E_I,\quad
 \sup_{N\in[X/2,X]\cap 2\Z}|e_X(N)|=o_w(X).
\end{equation}
\B{The original stage 75 statement was made on $H$. The inherited hypotheses use no primality condition on $N-1$, so the same uniform error applies to all even targets in this interval, in particular to $D$. At a prime midpoint the weighted count includes the diagonal; on $D$ it does not. In any of these domains the assertion $A_3=o_w(X)$ uniformly is equivalent to $r_X=F_X+o_w(X)$ uniformly. For a fixed $\lambda\ge0$, an eventual lower bound $A_3\ge-(\lambda+\epsilon)X$ for every $\epsilon>0$ is equivalent to the eventual bound $r_X\ge F_X-(\lambda+\epsilon)X$ with the same quantifiers: use $\epsilon/2$ in the premise and $|e_X|\le\epsilon X/2$ in either direction. This does not identify a finite threshold without slack. In particular, demanding $A_3=o_w(X)$ is demanding a weighted Hardy--Littlewood asymptotic, not an independent shortcut to it.}{El enunciado original de la etapa 75 se formuló en $H$. Las hipótesis heredadas no usan ninguna condición de primalidad de $N-1$, de modo que el mismo error uniforme se aplica a todos los pares del intervalo y, en particular, a $D$. Si el punto medio es primo, el recuento ponderado incluye la diagonal; en $D$ no. En cualquiera de estos dominios, $A_3=o_w(X)$ uniforme equivale a $r_X=F_X+o_w(X)$ uniforme. Para $\lambda\ge0$ fijo, una cota eventual $A_3\ge-(\lambda+\epsilon)X$ para todo $\epsilon>0$ equivale a la cota eventual $r_X\ge F_X-(\lambda+\epsilon)X$ con los mismos cuantificadores: se toma $\epsilon/2$ en la premisa y $|e_X|\le\epsilon X/2$ en ambos sentidos. Esto no identifica un umbral finito sin holgura. En particular, exigir $A_3=o_w(X)$ es exigir una asintótica ponderada de Hardy--Littlewood, no un atajo independiente hacia ella.}

\paragraph{\B{Conditional positivity.}{Positividad condicional.}}
\B{Choose one fixed broad weight $w>0$ on $[1/5,3/5]$ and put $m_w=\min_{[9/40,21/40]}w$. For $\tau\in[1/2,1]$, both $u$ and $\tau-u$ lie in this last interval when $u\in[\tau/2-1/40,\tau/2+1/40]$. Thus $J_w(\tau)\ge m_w^2/20$ and $F_X\ge c_w X$, with $c_w=C_2m_w^2/10>0$. The first missing theorem for this route is an independently justified bound on the signed Type II sum for the unchanged prime problem. One sufficient form is}{Elegimos un peso fijo de soporte amplio con $w>0$ en $[1/5,3/5]$ y ponemos $m_w=\min_{[9/40,21/40]}w$. Para $\tau\in[1/2,1]$, tanto $u$ como $\tau-u$ pertenecen a este último intervalo si $u\in[\tau/2-1/40,\tau/2+1/40]$. Así, $J_w(\tau)\ge m_w^2/20$ y $F_X\ge c_w X$, con $c_w=C_2m_w^2/10>0$. El primer teorema que falta en esta vía es una cota justificada de forma independiente para la suma con signo de Tipo II del problema original con primos. Una forma suficiente es}
\begin{equation}\label{eq:missing}
 \exists X_0\ \forall X\ge X_0\ \forall N\in D\cap[X/2,X]:
 \quad A_3(N)\ge-\tfrac12c_wX .
\end{equation}
\B{For the author’s statement alone replace $D$ by $H$. Once $|e_X|\le c_wX/4$, the proposed premise and \eqref{eq:transport} would give $r_X\ge c_wX/4>0$, hence $R(N)>0$ on the chosen domain. Taking $X=N$ proves the corresponding eventual assertion under this premise, not the universal assertion. The weight is important: if supported in $(1/10,1/5)$, then $r_X(N)=0$ on the entire target interval, regardless of primality.}{Para el enunciado del autor basta sustituir $D$ por $H$. Una vez que $|e_X|\le c_wX/4$, la premisa propuesta y \eqref{eq:transport} darían $r_X\ge c_wX/4>0$ y, por tanto, $R(N)>0$ en el dominio elegido. Tomar $X=N$ demuestra la afirmación eventual correspondiente bajo esa premisa, no la universal. El peso importa: si su soporte está en $(1/10,1/5)$, entonces $r_X(N)=0$ en todo el intervalo objetivo, con independencia de la primalidad.}

\paragraph{\B{What remains open.}{Qué queda abierto.}}
\B{The signed bound \eqref{eq:missing} is not proved. A universal conclusion by this route would also require a rigorous effective cutoff for the bound and errors and a finite verification covering the entire range below it. The reviewed bibliography and the two earlier papers supply no such input. The next admissible analytic step must identify information sensitive to the true reflected prime correlation beyond the data preserved in the deletion model. Repeating the wheel calculation, proving positivity of more restricted cores, or replacing a pointwise estimate by an exceptional-set bound cannot fill this gap. A valid logically different route to Goldbach is not excluded.}{La cota con signo \eqref{eq:missing} no está demostrada. Para una conclusión universal por esta vía se necesitaría además un umbral efectivo riguroso para la cota y los errores, y una verificación finita que cubra todo el tramo inferior. Ni la bibliografía consultada ni los dos artículos anteriores aportan ese ingrediente. El siguiente paso analítico admisible debe identificar información sensible a la correlación reflejada real, más allá de los datos que conserva el modelo de borrado. Repetir la rueda, probar la positividad de más núcleos restringidos o sustituir una cota puntual por una del conjunto excepcional no cubre esta laguna. No se excluye una vía válida lógicamente distinta para Goldbach.}

\subsection{\B{A two-sided estimate on \texorpdfstring{$H$}{H} and exceptional real zeros}{Una estimación bilateral en \texorpdfstring{$H$}{H} y los ceros reales excepcionales}}
\B{The count in this subsection is the uncut ordered von Mangoldt convolution $r_2(N)=\sum_{u+v=N}\Lambda(u)\Lambda(v)$, not the smooth genuine-prime count $r_X(N)$. We adapt the argument of Bhowmik and Grimmelt \cite[Proposition 8.1 and Theorem 8.2]{bg}: the additional step is to retain enough conductor multiples after imposing $N\in H$.}{El recuento de este apartado es la convolución ordenada de von Mangoldt sin corte $r_2(N)=\sum_{u+v=N}\Lambda(u)\Lambda(v)$, no el recuento suavizado de primos auténticos $r_X(N)$. Adaptamos el argumento de Bhowmik y Grimmelt \cite[Proposición 8.1 y Teorema 8.2]{bg}: el paso adicional es conservar suficientes múltiplos del conductor tras imponer $N\in H$.}

\noindent\textbf{\B{Proposition (conditional implication on $H$).}{Proposición (implicación condicional sobre $H$).}}
\emph{\B{Fix $A>5/2$, $0<\delta<1$, and any nonnegative function $\eta(r)\to0$. Put $X=r^A$. For each sufficiently large integer $r$, if}{Fijemos $A>5/2$, $0<\delta<1$ y una función no negativa $\eta(r)\to0$. Sea $X=r^A$. Para cada entero $r$ suficientemente grande, si}}
\begin{equation}\label{eq:siegelband}
 \delta\Sg(N)N\le r_2(N)\le(2-\delta)\Sg(N)N
\end{equation}
\emph{\B{holds at all $N\in H\cap[X/2,X]$ divisible by $r$, apart from at most $\eta(r)X/r$ exceptions, then no primitive real character of conductor $r$ has a real $L$-zero $\beta>1-c(\delta)/\log X$. The threshold for $r$ may depend on $A,\delta,\eta$.}{se cumple para los $N\in H\cap[X/2,X]$ divisibles por $r$, salvo a lo sumo $\eta(r)X/r$ excepciones, entonces ningún carácter real primitivo de conductor $r$ tiene un cero real $\beta>1-c(\delta)/\log X$ en su función $L$. El umbral para $r$ puede depender de $A,\delta,\eta$.}}

\begin{proof}
\B{Use only multiples of $2r$ in $[X/2,X]$. There are $X/(4r)+O(1)$, and $N/2$ is a multiple of $r$ strictly greater than $r$ for large $r$, hence composite. The only remaining way to fall outside $H$ is for $N-1$ to be prime. Such predecessors lie in the reduced class $-1$ modulo $r$. Brun--Titchmarsh and $r/\varphi(r)\ll\log\log r$ bound their number by}{Usamos solo los múltiplos de $2r$ en $[X/2,X]$. Hay $X/(4r)+O(1)$, y $N/2$ es un múltiplo de $r$ estrictamente mayor que $r$ para $r$ grande, luego compuesto. La única forma restante de quedar fuera de $H$ es que $N-1$ sea primo. Tales predecesores están en la clase reducida $-1$ módulo $r$. Brun--Titchmarsh y $r/\varphi(r)\ll\log\log r$ acotan su número por}
\[
 \frac{2X}{\varphi(r)\log(X/r)}\ll_A\frac{X}{r}\frac{\log\log r}{\log r}=o(X/r).
\]
\B{Suppose a primitive real $\chi$ of conductor $r$ has a real zero $\beta=1-\xi$ with $0<\xi\log X<c$. In \cite[Proposition 8.1]{bg} take $\vartheta=2/5$, so $r=X^{1/A}<X^{2/5}$, and choose $0<\epsilon<1-1/A-3/5$. The $O_\epsilon(X^{3/5+\epsilon})$ exceptions are $o(X/r)$. After excluding these, the prime predecessors, and the exceptions in \eqref{eq:siegelband}, at least one multiple remains. For it the generalized singular series in the exceptional-zero term equals $\chi(-1)\Sg(N)$; division by $\Sg(N)N$ gives}{Supongamos que un carácter real primitivo $\chi$ de conductor $r$ tiene un cero real $\beta=1-\xi$ con $0<\xi\log X<c$. En \cite[Proposición 8.1]{bg} tomamos $\vartheta=2/5$, de modo que $r=X^{1/A}<X^{2/5}$, y elegimos $0<\epsilon<1-1/A-3/5$. Las $O_\epsilon(X^{3/5+\epsilon})$ excepciones son $o(X/r)$. Excluidas estas, los predecesores primos y las excepciones de \eqref{eq:siegelband}, queda al menos un múltiplo. En él la serie singular generalizada del término de cero excepcional vale $\chi(-1)\Sg(N)$; al dividir por $\Sg(N)N$ se obtiene}
\begin{align*}
 \frac{r_2(N)}{\Sg(N)N}&=1+\chi(-1)\,\mathrm B(\beta,\beta)N^{-2\xi}+O\!\left(c^{c'}+\frac{r}{\varphi(r)^2}+X^{-\epsilon_0}\right).
\end{align*}
\B{Here $\mathrm B$ is the beta function and the last bound uses $X/2\le N\le X$ and $\Sg(N)\ge2C_2$. Uniformly on this interval, the beta factor times $N^{-2\xi}$ is $1+O(c)$, by smoothness near $(1,1)$ and $\xi\log N\le c$. Choosing $c$ small and then $r$ large makes the normalized count lie within strictly less than $\delta$ of $0$ if $\chi(-1)=-1$, or of $2$ if $\chi(-1)=1$. Either case contradicts \eqref{eq:siegelband}.}{Aquí $\mathrm B$ es la función beta y la última cota usa $X/2\le N\le X$ y $\Sg(N)\ge2C_2$. Uniformemente en este intervalo, el factor beta multiplicado por $N^{-2\xi}$ es $1+O(c)$, por suavidad cerca de $(1,1)$ y $\xi\log N\le c$. Eligiendo primero $c$ pequeño y después $r$ grande, el recuento normalizado queda a distancia estrictamente menor que $\delta$ de $0$ si $\chi(-1)=-1$, o de $2$ si $\chi(-1)=1$. En ambos casos se contradice \eqref{eq:siegelband}.}
\end{proof}

\B{The band is unproved; the conclusion is a real zero-free region at the indicated scale, not a fixed zero-free half-plane. A uniform uncut asymptotic $r_2(N)\sim\Sg(N)N$ on $H$ would give the band. Neither mere existence of a prime pair nor the one-sided Type II margin gives it. In particular $A_3=o_w(X)$ for one fixed smooth weight yields a weighted asymptotic; relating that to this uncut two-sided estimate would require an additional argument. This analytic consequence is not an obstruction that every proof of Goldbach must overcome.}{La banda no está demostrada; la conclusión es una región libre de ceros reales a la escala indicada, no un semiplano fijo libre de ceros. Una asintótica uniforme sin corte $r_2(N)\sim\Sg(N)N$ en $H$ sí daría la banda. Ni la mera existencia de una pareja prima ni el margen unilateral de Tipo II la proporcionan. En particular, $A_3=o_w(X)$ para un solo peso suave fijo da una asintótica ponderada; relacionarla con esta estimación bilateral sin corte requeriría otro argumento. Esta consecuencia analítica no es una obstrucción que deba superar toda prueba de Goldbach.}

\section{\B{A check on a proposed antidiagonal argument}{Comprobación de un argumento por antidiagonales}}
\B{Qazaqzeh and Rawashdeh \cite[Definition 2.1]{qr} use their symbol $\mathcal O(g)$ for convergence of the ratio to a nonzero constant, rather than Landau's upper-bound convention. Their argument on p.~5 asserts $R_1=R_2$ and $R_3=R_4$ by symmetry. Both equalities fail. At $2n=100$, on the common domain $x,y-2\ge3$ with $x,y$ odd and $x+y=100$, the four patterns from their lists give the following values. The symbols $P,C$ denote prime and composite, respectively.}{Qazaqzeh y Rawashdeh \cite[Definición 2.1]{qr} usan su símbolo $\mathcal O(g)$ para la convergencia del cociente a una constante no nula, no para la cota superior de Landau. Su argumento de la p.~5 afirma $R_1=R_2$ y $R_3=R_4$ por simetría. Ambas igualdades fallan. Para $2n=100$, en el dominio común $x,y-2\ge3$ con $x,y$ impares y $x+y=100$, los cuatro patrones de sus listas dan los siguientes valores. Los símbolos $P,C$ denotan primo y compuesto, respectivamente.}
\begin{center}\small
\begin{tabular}{@{}lccc@{}}
\toprule
\B{Count}{Recuento}&$(x,x+2,y,y-2)$&\B{Eligible $x$}{Valores de $x$}&\B{Total}{Total}\\
\midrule
$R_1$&$P,P,C,P$&---&0\\
$R_2$&$C,P,C,C$&35, 65&2\\
$R_3$&$P,P,P,C$&3, 11, 17, 29, 41, 59, 71&7\\
$R_4$&$P,C,C,C$&7, 13, 23, 43, 73&5\\
\bottomrule
\end{tabular}
\end{center}
\B{All entries can be checked by trial division; the certificate accompanies the article. In fact the mismatch persists for large even $2n$: $R_1,R_3=O(n/\log^2 n)$ by the upper sieve for twin primes \cite{wu}, while $R_2,R_4\sim 2n/\log(2n)$. For the latter, start from the choices with the required single prime ($x+2$ or $x$) and exclude cases in which one of the other three entries is also prime. Each exclusion is either a gap-two prime pair or a pair of sum $2n$ or $2n\pm2$. The upper Goldbach sieve bounds each sum constraint by $O(n\log\log n/\log^2 n)$ uniformly; boundary exclusions cost $O(1)$. The argument therefore cannot cancel the single-prime configurations by the claimed equalities. This rejects the proof step, not the Goldbach conjecture.}{Todas las entradas se comprueban por división por tentativa; el certificado acompaña al artículo. La discrepancia persiste para $2n$ par grande: $R_1,R_3=O(n/\log^2 n)$ por la criba superior de primos gemelos \cite{wu}, mientras que $R_2,R_4\sim 2n/\log(2n)$. Para estos últimos se parte de las opciones con el único primo exigido ($x+2$ o $x$) y se retiran las que también tienen otro primo entre las tres entradas restantes. Cada exclusión impone una pareja de diferencia dos o de suma $2n$ o $2n\pm2$. La criba superior de Goldbach acota cada condición de suma uniformemente por $O(n\log\log n/\log^2 n)$; los bordes cuestan $O(1)$. Por tanto, las igualdades alegadas no permiten cancelar las configuraciones de un solo primo. Se invalida el paso de la prueba, no la conjetura de Goldbach.}

\paragraph{\B{An independent test of the constant-limit claim.}{Una comprobación independiente del supuesto límite constante.}}
\B{Let $G(N)$ count ordered odd-prime pairs. The ratio $G(N)/(N/\log^2 N)$ cannot converge to a finite constant along all even $N$. By the prime number theorem in arithmetic progressions, for $a=1,5$ the measures $(2\log T/T)\sum_{p\le T,\,p\equiv a\,(6)}\delta_{p/T}$ converge weakly to Lebesgue measure on $[0,1]$. Products converge on the triangle $u+v\le1$ (its boundary has measure zero). Thus each ordered pair of reduced classes contributes asymptotically $T^2/(8\log^2 T)$ pairs of sum at most $T$. For $N\equiv0\pmod6$ there are two such ordered class pairs, and for $N\equiv2\pmod6$ only one; pairs involving 3 contribute $O(T)$. If the ratio converged to $L$, the two sums of $G(N)$ over these classes would both equal $(L+o(1))T^2/(12\log^2 T)$, by summing $N/\log^2 N$ on each progression. This contradicts the unequal coefficients $1/4$ and $1/8$. No conjectural asymptotic for individual targets is used. Replacing the denominator by $n/\log^2 n$ at $N=2n$ changes only its asymptotic constant factor.}{Sea $G(N)$ el recuento ordenado de parejas de primos impares. El cociente $G(N)/(N/\log^2 N)$ no puede converger a una constante finita a lo largo de todos los pares $N$. Por el teorema de los números primos en progresiones, para $a=1,5$ las medidas $(2\log T/T)\sum_{p\le T,\,p\equiv a\,(6)}\delta_{p/T}$ convergen débilmente a la medida de Lebesgue en $[0,1]$. Sus productos convergen en el triángulo $u+v\le1$ (la frontera tiene medida nula). Así, cada pareja ordenada de clases reducidas aporta asintóticamente $T^2/(8\log^2 T)$ parejas de suma hasta $T$. Para $N\equiv0\pmod6$ hay dos parejas ordenadas de clases y para $N\equiv2\pmod6$ solo una; las parejas con el 3 aportan $O(T)$. Si el cociente convergiese a $L$, ambas sumas de $G(N)$ por clase serían $(L+o(1))T^2/(12\log^2 T)$, al sumar $N/\log^2 N$ en cada progresión. Esto contradice los coeficientes distintos $1/4$ y $1/8$. No se usa ninguna asintótica conjetural por objetivo. Sustituir el denominador por $n/\log^2 n$ cuando $N=2n$ solo cambia su factor constante asintótico.}


\section{\B{Evidence record and reproducibility}{Registro de evidencia y reproducibilidad}}
\B{The project record uses the selected Stage 76 snapshot (PR 27, commit \texttt{d9e0ac05\allowbreak 247f2f22\allowbreak 0423c444\allowbreak d3056164\allowbreak 5c747c4c}); Stage 73's proof ledger and the Stage 74 and 75 deltas supply the theorem identifiers. The governance dossier (PR 28) records the separate review scopes. These are internal provenance references, not publicly accessible evidence required to read the paper: the principal derivations and their qualifications are printed above. Early project trials (stages 56--69) are consolidated by mechanism rather than recounted as independent advances. Fixed-wheel counts and matching fixtures in the historical record were checked in Python and Sage; the present experiment uses only the separately supplied Python source. No new Lean compilation or formalization of the pair theorem is represented.}{El registro del proyecto utiliza la versión elegida de la etapa 76 (PR 27, commit \texttt{d9e0ac05\allowbreak 247f2f22\allowbreak 0423c444\allowbreak d3056164\allowbreak 5c747c4c}); el registro de la etapa 73 y las adendas de las etapas 74 y 75 aportan los identificadores. El expediente organizativo (PR 28) registra los ámbitos separados de revisión. Son referencias internas de procedencia, no pruebas públicas necesarias para leer el artículo: las derivaciones principales y sus salvedades se exponen arriba. Los ensayos tempranos (etapas 56--69) se sintetizan por mecanismo y no se presentan como avances independientes. Los recuentos de la rueda y los casos de emparejamiento del registro histórico se comprobaron con Python y Sage; el experimento actual emplea solo el código Python adjunto. No se afirma una nueva compilación de Lean ni una formalización del teorema de parejas.}

\B{Evidence labels: the earlier Lean results are reported as prior results; the logical and finite-sieve identities are elementary derivations with explicit proofs and separate finite checks; the table is a computational observation; T73.5/T74.3/T75.1 are internal derivations pending independent analytical review; the Type II implication is conditional; the positivity inference from deletion-stable data in the enlarged class is refuted. Both universal assertions remain open. To reproduce the new table, run \texttt{python3 goldbach\_reflected\_audit.py 7930} with the standard library and compare with the supplied JSON.}{Etiquetas: los resultados en Lean se presentan como antecedentes; las identidades lógicas y de criba son derivaciones elementales con pruebas explícitas y comprobaciones finitas separadas; la tabla es una observación computacional; T73.5/T74.3/T75.1 son derivaciones internas pendientes de revisión analítica independiente; la implicación de Tipo II es condicional; se refuta la inferencia de positividad a partir de datos estables tras el borrado en la clase ampliada. Ambas afirmaciones universales siguen abiertas. Para reproducir la tabla, ejecútese \texttt{python3 goldbach\_reflected\_audit.py 7930} con la biblioteca estándar y compárese con el JSON adjunto.}

\B{The extended witness audit and the antidiagonal certificate can be regenerated with the C++17 and Python programs in the companion archive. The real-zero implication is a deduction from an external formula, with an explicit unproved premise, and has not received independent referee review. Neither it nor the fixed-modulus average establishes a new target-pointwise lower bound.}{La comprobación ampliada de testigos y el certificado de antidiagonales se reproducen con los programas en C++17 y Python del archivo complementario. La implicación sobre ceros reales es una deducción de una fórmula externa, con una premisa no demostrada explícita, y no ha recibido revisión independiente. Ni ella ni el promedio en módulos fijos establecen una nueva cota inferior puntual.}

\begin{thebibliography}{99}\small\interlinepenalty=10000
\bibitem{additive} T. Sánchez Expósito, \emph{An Additive Sieve on an Arithmetic Progression for Prime Number Generation: Formal Verification in Lean 4}, SSRN 5705571, revised 2026. \href{https://doi.org/10.2139/ssrn.5705571}{doi:10.2139/ssrn.5705571}.
\bibitem{twins} T. Sánchez Expósito, \emph{Compressed Residue Sieve for Twin-Prime Candidate Verification: Formal Verification in Lean 4}, \href{https://papers.ssrn.com/sol3/papers.cfm?abstract_id=7489199}{SSRN 7489199}, 2026. Software: \href{https://doi.org/10.5281/zenodo.22846302}{doi:10.5281/zenodo.22846302}.
\bibitem{ford} K. Ford, \emph{Sieve Methods Lecture Notes}, 2023, §§1.2 and 2. \url{https://ford126.web.illinois.edu/sieve2023.pdf}.
\bibitem{verify} T. Oliveira e Silva, S. Herzog and S. Pardi, Empirical verification of the even Goldbach conjecture and computation of prime gaps up to $4\cdot10^{18}$, \emph{Math. Comp.} 83 (2014), 2033--2060. \href{https://doi.org/10.1090/S0025-5718-2013-02787-1}{doi:10.1090/S0025-5718-2013-02787-1}.
\bibitem{chen} J.-R. Chen, On the representation of a larger even integer as the sum of a prime and the product of at most two primes, \emph{Sci. Sinica} 16 (1973), 157--176.
\bibitem{pintz} J. Pintz, A new explicit formula in the additive theory of primes with applications II. The exceptional set in Goldbach's problem, 2018, \href{https://arxiv.org/abs/1804.09084}{arXiv:1804.09084}.
\bibitem{helfgott} H. A. Helfgott, The ternary Goldbach conjecture is true, 2013, \href{https://arxiv.org/abs/1312.7748}{arXiv:1312.7748}.
\bibitem{mv} H. L. Montgomery and R. C. Vaughan, \emph{Multiplicative Number Theory I: Classical Theory}, Cambridge University Press, 2007, Ch.~6.
\bibitem{bv} E. W. Weisstein, Bombieri--Vinogradov Theorem, \emph{MathWorld} (maximal discrepancy formulation). \url{https://mathworld.wolfram.com/Bombieri-VinogradovTheorem.html}.
\bibitem{wu} J. Wu, Chen's double sieve, Goldbach's conjecture and the twin prime problem, \emph{Acta Arith.} 114 (2004), 215--273. \href{https://doi.org/10.4064/aa114-3-2}{doi:10.4064/aa114-3-2}.

\bibitem{oneill} M. E. O'Neill, The genuine sieve of Eratosthenes, \emph{J. Funct. Program.} 19 (2009), 95--106. \href{https://doi.org/10.1017/S0956796808007004}{doi:10.1017/S0956796808007004}.
\bibitem{zhao} G. Zhao, The exceptional set of Goldbach problem and Linnik's constant, preprint, January 2026 revision, \href{https://arxiv.org/abs/2511.05631v2}{arXiv:2511.05631v2}.
\bibitem{bg} G. Bhowmik and L. Grimmelt, The exceptional set of the Goldbach problem, \emph{Rend. Circ. Mat. Palermo} (2026), §8; \href{https://arxiv.org/abs/2607.27282v1}{arXiv:2607.27282v1}; \href{https://doi.org/10.1007/s10476-026-00191-z}{doi:10.1007/s10476-026-00191-z}.
\bibitem{pa} J. Pintz and I. Althöfer, The Goldbach exceptional set: an exponent 7/12 and conditional refinement, preprint, 2026, \href{https://arxiv.org/abs/2610.15559}{arXiv:2610.15559}.
\bibitem{mpr} J. Maynard, M. Pandey and M. Radziwiłł, Exponential sums over primes, preprint, 2026, \href{https://arxiv.org/abs/2608.14777}{arXiv:2608.14777}.
\bibitem{qr} K. Qazaqzeh and A. Rawashdeh, An elementary proof of the binary Goldbach conjecture, preprint, 6 November 2026, \href{https://arxiv.org/abs/2611.09007v1}{arXiv:2611.09007v1}. Cited for the claim examined here.


\end{thebibliography}
\end{document}
